% ------------------------------------------------------------------------------
% Plantilla de tesis - Universidad de Antioquia
% Copyright (c) 2026 Simón Zuluaga Henao
%
% Este archivo hace parte de una adaptación de la "Plantilla (LaTeX) IEEE"
% del Sistema de Bibliotecas de la Universidad de Antioquia.
%
% Distribuido bajo la licencia CC BY-NC-SA 4.0.
% ------------------------------------------------------------------------------

% ==================================================
% %%%%% ESTE ARCHIVO NO DEBE SER EDITADO %%%%%
% ==================================================



% ==================================================
% CONFIGURACIÓN
% ==================================================

\usepackage[top=2.5cm, bottom=3cm, left=3cm, right=3cm]{geometry}
\usepackage[table]{xcolor}
\usepackage{graphicx,amsmath,amssymb,array,enumitem,lipsum,float,etoolbox,xstring}



\usepackage{caption}
\newcommand{\textFigure}{Fig.}
\DeclareCaptionTextFormat{up}{\MakeUppercase{#1}}
\captionsetup[table]{labelformat={default}, font=footnotesize, justification=centering, labelsep=newline, textformat=up}
\renewcommand{\thetable}{\Roman{table}}
\captionsetup[figure]{labelformat={default}, font=footnotesize, labelsep=period, justification=centering,singlelinecheck=false, name={\textFigure}}
\renewcommand{\thefigure}{\arabic{figure}}



\usepackage[numbers]{natbib}
\setlength{\parindent}{1.25cm}          % Sangria
\renewcommand{\baselinestretch}{1.5}    % Interlineado
\setlength{\parskip}{0.1em}             % Espacio entre parrafos
\bibliographystyle{ieeeudea}

\renewcommand{\labelitemi}{$\bullet $}
\setlist[itemize]{leftmargin=*, parsep=-5px}
\renewcommand\descriptionlabel[1]{$\bullet$\hspace{0.5em} \textbf{#1}}
\setlist[description]{parsep=-5px}



\usepackage{titletoc}
\titlecontents{section}
  [0pt]
  {\smallskip}
  {\makebox[2.5em][l]{\hspace{4pt}\thecontentslabel.}}
  {}
  {\titlerule*[1pc]{.}\contentspage}
\titlecontents{subsection}
  [1.5em]
  {\smallskip}
  {\makebox[2.5em][l]{\thecontentslabel.}}
  {}
  {\titlerule*[1pc]{.}\contentspage}
\titlecontents{subsubsection}
  [3.5em]
  {\smallskip}
  {\makebox[2.5em][l]{\thecontentslabel)}}
  {}
  {\titlerule*[1pc]{.}\contentspage}

\setcounter{tocdepth}{3}
\setcounter{secnumdepth}{4}     % Define el contador para los 4 niveles de titulo

\renewcommand{\thesection}{\Roman{section}}
\renewcommand{\thesubsection}{\Alph{subsection}}
\renewcommand{\thesubsubsection}{\arabic{subsubsection}}
\renewcommand{\theparagraph}{\alph{paragraph}}



\usepackage{titlesec}
\titleformat{\section}{\centering\normalfont\normalsize\MakeUppercase}{\thesection.}{0.2em}{}
\titleformat{\subsection}{\normalfont\normalsize\it}{\thesubsection.}{0.2em}{}
\titleformat{\subsubsection}[runin]{\itshape\normalsize\it}{\hspace{0.5cm}\thesubsubsection)}{0.2em}{}[:]



\usepackage{fancyhdr}
\definecolor{verdeUdeA}{rgb}{0.325,0.506,0.208}     % RGB=83,129,53
\setlength{\headheight}{25pt}
\renewcommand{\headrule}{{\color{verdeUdeA}\rule{\textwidth}{1pt}}}
\renewcommand{\headrulewidth}{2pt}
\lhead{\footnotesize\MakeUppercase{\tituloCorto}}\chead{}\rhead{{\color{white}Rufino}\thepage}
\cfoot{}



\usepackage{hyperref}
% \hypersetup{colorlinks, linkcolor=blue, urlcolor=blue, citecolor={red}}       % Color links
\hypersetup{hidelinks}          % Black links



\usepackage{booktabs}
\setlength{\tabcolsep}{8pt}         % Horizontal padding



% LOT, LOF configuration
\makeatletter
\let\oldl@table\l@table
\renewcommand*\l@table[2]{%
  \begingroup
  \renewcommand{\numberline}[1]{\textTable~##1\hspace{1em}}%
  \oldl@table{#1}{#2}%
  \endgroup
}
\makeatother

\makeatletter
\let\oldl@figure\l@figure
\renewcommand*\l@figure[2]{%
  \begingroup
  \renewcommand{\numberline}[1]{\textFigure~##1\hspace{1em}}%
  \oldl@figure{#1}{#2}%
  \endgroup
}
\makeatother



% Bookmarks setup
\newcommand{\tocuppercase}[1]{
  \texorpdfstring{\MakeUppercase{#1}}{#1}
}
\let\oldsection\section
\RenewDocumentCommand{\section}{s o m}{
  \IfBooleanTF{#1}
    {% \section*
      \oldsection*{#3}%
    }
    {% \section
      \IfNoValueTF{#2}
        {\oldsection[\tocuppercase{#3}]{#3}}
        {\oldsection[#2]{#3}}
    }
}
\newcommand{\sectionCustom}[1]{
  \section*{#1}
  \addcontentsline{toc}{section}{\tocuppercase{#1}}
}
